\section{Ingeniería y monitoreo}

\subsection{RLHF Pipeline Timeline}

Dividir el RLHF pipeline en etapas facilita la colaboración del equipo:

\begin{itemize}
    \item Semana de construcción de datos: recolección de instruction + comparison data
    \item Semana de modelado de recompensa: entrenamiento repetido + calibration + validation
    \item Semana de optimización de la política: iteraciones de PPO o DPO + KL/entropy tuning
    \item Semana de control de calidad: benchmark + red team + human eval
    \item Monitoreo y rollout: gradual deployment + drift monitoring
\end{itemize}

\begin{importantbox}{Action Timeline}
Cada etapa define un deliverable claro: resultados de la construcción de datos, blueprints del modelo de recompensa, checkpoint de PPO, informe de revisión y decisión de despliegue. Solo así el equipo sabe cuándo puede avanzar al siguiente paso.
\end{importantbox}

\subsection{Operational Monitoring}

Después del despliegue es indispensable monitorear de forma continua hallucination, rejection y latency.

\begin{figure}[H]
\centering
\includegraphics[width=0.8\textwidth]{slide-01.jpg}
\caption{La Slide 1 muestra a la vez el RLHF pipeline y las métricas de monitoreo, y enfatiza el seguimiento de latency y rejection.}
\end{figure}

\begin{knowledgebox}{Diseño de las métricas del tablero de observación}
Se recomienda mostrar al mismo tiempo: Reward score drift, human rejection rate, token latency y deploy coverage. Combinar distintas métricas ayuda al equipo a determinar con rapidez si se trata de model drift o de data drift.
\end{knowledgebox}

\begin{table}[H]
\centering
\begin{tabular}{p{0.3\textwidth} p{0.6\textwidth}}
\toprule
Métrica & Descripción \\
\midrule
Reward score drift & Compara con una rolling window la distribución de recompensa current vs baseline \\
Reject ratio & Tasa de rechazo de la retroalimentación de revisión; sirve para evaluar la sensibilidad de los guardrails \\
Latency P99 & Garantiza que el pipeline de transformer + verifier no exceda el SLO \\
Human override count & Número de intervenciones humanas; sirve para calibrar la automation trust \\
\bottomrule
\end{tabular}
\caption{Métricas centrales del tablero de observación posterior al despliegue}
\end{table}

\subsection{Governance Checklist}

\begin{table}[H]
\centering
\begin{tabular}{p{0.2\textwidth} p{0.35\textwidth} p{0.28\textwidth}}
\toprule
Punto de control & Pregunta central & Salida \\
\midrule
Reward model approval & ¿Pasó por blind eval validation? & Approval checklist \\
Policy guardrails & ¿Existe automatic rejection for forbidden topics? & Guardrail rules \\
Deployment gating & ¿Se amplía la capacidad de forma gradual y se retienen las anomalías? & Rollout policy \\
\bottomrule
\end{tabular}
\caption{Lista de verificación de gobernanza de RLHF}
\end{table}

\subsection{Arquitectura de Drift Detection}

Una vez que el modelo está en producción, hay que detectar dos tipos de drift: prompt drift (cambios en la entrada del usuario) y distribution drift (cambios en el comportamiento de la salida).

\begin{itemize}
    \item Capturar en el logging pipeline la prompt signature (hash + metadata)
    \item Calcular con una rolling window la distribución de reward score
    \item Proporcionar una alert que active una human review cuando KL divergence vs zscore supere el umbral
    \item Crear un `drift ticket` que incluya snapshot prompt, reward score y output diff
\end{itemize}

\begin{knowledgebox}{Drift detection no es solo anomaly detection}
Requiere context-aware thresholds: distintas tareas (summaries vs math) tienen de por sí una variance distinta en el reward score. Sin un context-specific threshold, se generan muchas falsas alarmas.
\end{knowledgebox}

\subsection{Resumen de la sección}

En el plano de la ingeniería hay que integrar el timeline, el monitoreo y la gobernanza; el dashboard de monitoreo y la checklist son una base importante de la confianza posterior al despliegue.
\subsection{Ops Playbook}

\begin{enumerate}
    \item \textbf{Observe}: agregar reward scores, rejection logs y prompt hash.
    \item \textbf{Assess}: revisar el drift con una evidence matrix y counterfactual prompts.
    \item \textbf{Act}: hacer que la IA recomiende acciones de mitigación low-risk y genere un rollback plan, que se ejecuta después de la aprobación humana.
    \item \textbf{Document}: registrar la decisión en la checklist / runbook para usarla en la siguiente ronda de evaluation.
\end{enumerate}

\begin{table}[H]
\centering
\begin{tabular}{p{0.18\textwidth} p{0.4\textwidth} p{0.3\textwidth}}
\toprule
Etapa & Salida & Participantes principales \\
\midrule
Observe & Reward + guardrail logs & SRE + Data Eng \\
Assess & Evidence scorecard + drift report & Alignment + QA \\
Act & Rollback plan/mitigation action & Product + Legal \\
Document & Runbook entry + retrospective & Ops + PM \\
\bottomrule
\end{tabular}
\caption{Entregables ejecutables del Ops Playbook}
\end{table}

\begin{knowledgebox}{El valor del Ops runbook}
Registrar cada ciclo de Act como una entrada del runbook no solo permite guarantee la reproducibilidad, sino que también ofrece un audit trail transparente en la board review.
\end{knowledgebox}

\begin{warningbox}{La IA se encarga principalmente del triage; la decisión sigue requiriendo supervisión humana}
Aunque la IA sugiera una regla, en las líneas de alto riesgo una persona debe revisarla, para evitar que un error de juicio de la automatización tenga consecuencias irreversibles.
\end{warningbox}

\subsection{Resumen de la sección}

En el plano de la ingeniería hay que encadenar el timeline, el monitoreo, la gobernanza, la drift detection y el ops playbook en un ciclo cerrado; solo así se mantiene la confianza en el alignment durante las etapas de despliegue gradual y despliegue completo.
